\section{Related Work}\label{sec:related}

\subsection{3D-Aware Vision-Language-Action Models}

A growing body of work integrates 3D spatial information into VLA architectures to improve manipulation performance. DepthVLA~\cite{depthvla2025} and QDepth-VLA~\cite{li2025qdepth} augment VLA training with depth prediction as an auxiliary objective, while Spatial Forcing~\cite{li2025spatialforcing} enforces implicit 3D alignment through spatial consistency constraints on the learned representations. Several methods go further by incorporating explicit 3D structure: 3D-CAVLA~\cite{3dcavla2025} and GeoAware-VLA~\cite{geoaware2025} condition action generation on point cloud features, SpatialVLA~\cite{spatialvla2025} encodes spatial affordance maps, and GST-VLA~\cite{gstvla2026}, StemVLA~\cite{stemvla2026}, and ST-VLA~\cite{stvla2026} fuse spatiotemporal representations from multi-view or sequential observations. Evo-0~\cite{evo02025}, SA-VLA~\cite{savla2026}, AugVLA-3D~\cite{augvla3d2026}, and Evo-Depth~\cite{evodepth2026} further confirm that 3D-aware features improve task success rate, viewpoint generalization, and grasping precision.

Among these, QDepth-VLA~\cite{li2025qdepth} is the closest methodological precedent for auxiliary spatial supervision in VLAs. Three distinctions separate it from our approach. First, QDepth-VLA predicts 2.5D camera-frame depth, which encodes per-pixel distance to the camera sensor; ColliForce-VLA predicts workspace-frame SDF, which encodes signed distance to the nearest obstacle surface at arbitrary 3D points. Second, QDepth-VLA's VQ-VAE decoder is architecturally decoupled from the action head: depth gradients flow into the backbone during training but do not explicitly modify the action token generation pathway. In contrast, our SDF gradient refinement directly modulates the feature representations from which action tokens are decoded. Third, camera-frame depth provides no signal in occluded regions (behind the robot arm or other objects), whereas workspace SDF covers the full 3D volume including occluded space. Spatial Forcing~\cite{li2025spatialforcing} demonstrates that implicit 3D alignment within the VLA can improve manipulation, yet it does not produce an occupancy or distance field and therefore cannot provide collision-specific directional guidance.

Despite these advances, no 3D-aware VLA targets collision avoidance as a design objective, and none evaluates on collision-specific metrics. ColliForce-VLA bridges this gap by using 3D geometric information (workspace SDF) explicitly for collision-aware action generation.

\subsection{Collision-Safe Robot Policies}

Recent methods address collision safety in learned manipulation policies through two strategies: post-hoc correction and constrained optimization. AEGIS~\cite{hu2025vlsa} exemplifies the post-hoc approach by applying Control Barrier Function (CBF) corrections to VLA-generated actions, using open-set detection and monocular depth estimation to identify obstacles at test time. CoFreeVLA~\cite{cofreeevla2026} trains a risk estimator that halts execution when predicted collision probability exceeds a threshold in dual-arm settings. On the constrained optimization side, SafeVLA~\cite{zhang2025safevla} formulates safety as a Constrained Markov Decision Process and reduces violations through reward shaping during reinforcement learning. Beyond VLA-specific methods, the broader robotics literature has explored collision-safe policies through inverse-kinematics-based quadratic programming with RL, cost function denoising conditioned on language instructions, and contact-aware policy execution.

These approaches share a common architectural pattern: the safety mechanism operates independently of the VLA's learned feature representations. The model generates actions without collision awareness, and a separate module evaluates or corrects them post-hoc, or a separate reward signal implicitly penalizes unsafe behavior without encoding geometric structure. ColliForce-VLA departs from this paradigm by internalizing collision geometry into the VLA's own parameters. The SDF gradient refinement during training steers the action expert's feature space toward collision-free configurations, so that the model encodes collision awareness at the point of action generation rather than relying on external correction.

\subsection{Signed Distance Fields in Robotics}

Signed Distance Fields provide a continuous implicit representation of geometry where each point encodes its signed distance to the nearest surface. DeepSDF~\cite{park2019deepsdf} established neural network-based SDF prediction as a compact shape representation, and iSDF~\cite{ortiz2022isdf} extended this to incremental scene reconstruction from posed depth images. DiffCo~\cite{zhi2021diffco} introduced differentiable collision checking based on learned distance fields, which enables gradient-based motion optimization through collision constraints. Subsequent work has applied SDFs broadly in robotics: RAMP and CSSDF-Net use SDF representations for real-time motion planning, Neural NMPC-SDF integrates SDF constraints into nonlinear model predictive control, and SAMP and Prof.\ Robot incorporate distance fields into policy learning for contact-rich tasks.

Two observations motivate our design. First, the effectiveness of SDF for collision detection in classical planning is well established, and DiffCo~\cite{zhi2021diffco} specifically demonstrates that differentiable SDF enables gradient-based collision avoidance. Second, existing SDF-based policy methods (SAMP, Prof.\ Robot) operate in non-VLA settings and do not handle language-conditioned manipulation. ColliForce-VLA is, to our knowledge, the first method to integrate a differentiable SDF into a VLA's forward pass, connecting the proven collision-detection capability of distance fields with the generalization capacity of large-scale vision-language-action models.

\subsection{Action Refinement in Generative Policies}

DFM-VLA has shown that iterative action refinement within the VLA architecture can improve action quality by denoising predicted trajectories through multiple refinement steps. This validates the feasibility of in-model refinement but lacks a collision-specific guidance signal; the refinement is driven by a general denoising objective rather than geometric constraints. ColliForce-VLA's SDF gradient refinement provides a collision-specific instantiation of this principle: the refinement direction is determined by the SDF gradient in feature space ($\nabla_z \Phi$), which explicitly encodes the direction of increasing distance from obstacle surfaces.
